\definecolor{dsdlcodebg}{RGB}{232,238,245}
\definecolor{dsdlcodeborder}{RGB}{208,215,222}
\newsavebox{\dsdlcodebox}
\newenvironment{dsdlcodeblock}
{%
  \par\noindent
  \begin{latin}
  \begin{lrbox}{\dsdlcodebox}
  \begin{minipage}{1.1\linewidth}
  \small
}
{%
  \end{minipage}
  \end{lrbox}
  \setlength{\fboxsep}{6pt}
  \fcolorbox{dsdlcodeborder}{dsdlcodebg}{\usebox{\dsdlcodebox}}
  \end{latin}
  \par
}

% -------------------------------------------------------
% Abstract
% -------------------------------------------------------

\بدون‌تورفتگی \مهم{چکیده}: در این آزمایش، یک سامانه‌ی \لر{UART} هفت‌بیتی با زبان توصیف سخت‌افزار
\لر{Verilog} طراحی و شبیه‌سازی شد. قالب فریم شامل ده بازه‌ی بیتی با ترتیب اختصاصی بیت شروع،
بیت توازن، هفت بیت داده از کم‌ارزش‌ترین بیت تا پرارزش‌ترین بیت و بیت توقف است. برای نرخ
۱۱۵۲۰۰ بیت بر ثانیه و ساعت ۵۰ مگاهرتز، هر بیت در ۴۳۴ چرخه‌ی ساعت نگه داشته می‌شود.
فرستنده و گیرنده به‌صورت دو ماشین حالت متناهی مستقل طراحی شدند و در پیمانه‌ی
\کد{UARTTop.v}، خروجی سریال فرستنده مستقیماً به ورودی گیرنده متصل شد. گیرنده با
هم‌زمان‌ساز دو مرحله‌ای و نمونه‌برداری در میانه‌ی هر بیت، خطاهای توازن، بیت توقف نامعتبر و
شروع کاذب را تشخیص می‌دهد. یازده آزمون خودکنترل اجرا شدند؛ همه‌ی آزمون‌ها، از جمله بررسی
تمام ۱۲۸ مقدار هفت‌بیتی، با موفقیت پایان یافتند. بررسی سنتز با \لر{Yosys} نیز صفر خطای
ساختاری و ۱۰۸ سلول در سلسله‌مراتب کامل طرح گزارش کرد.

\بدون‌تورفتگی \مهم{واژه‌های کلیدی}:
ارتباط سریال، بیت توازن، ماشین حالت متناهی، نمونه‌برداری میانی، آزمون‌گر خودکنترل،
\لر{UART}، \لر{Verilog}

\section{مقدمه و اهداف آزمایش}

\لر{UART}\پاورقی{\لر{Universal Asynchronous Receiver/Transmitter}} یک روش ارتباط سریال
ناهمگام است که داده‌ی موازی را در فرستنده به دنباله‌ای از بیت‌ها تبدیل و در گیرنده دوباره
بازسازی می‌کند. در این آزمایش، برخلاف قالب متداول \لر{UART}، بیت توازن پیش از بیت‌های داده
قرار می‌گیرد. ترتیب دقیق فریم عبارت است از:

\begin{dsdlcodeblock}
\begin{verbatim}
Start(0) -> Parity -> D0 -> D1 -> D2 -> D3 -> D4 -> D5 -> D6 -> Stop(1)
\end{verbatim}
\end{dsdlcodeblock}

هدف‌های اصلی آزمایش، طراحی فرستنده‌ی هفت‌بیتی، طراحی گیرنده با کنترل توازن و بیت توقف،
اتصال مستقیم دو پیمانه در حالت حلقه‌ی بازگشت\پاورقی{\لر{Loopback}}، پیاده‌سازی کد قابل سنتز
و ارزیابی رفتار طرح در حالت معتبر و حالت‌های خطا است.

\section{قالب فریم و محاسبات زمانی}

خط سریال در حالت بیکار مقدار منطقی یک دارد. شروع فریم با صفرشدن خط مشخص می‌شود. پس از
بیت شروع، بیت توازن و سپس هفت بیت داده به‌ترتیب \لر{LSB-first} ارسال می‌شوند. بیت توقف
مقدار یک دارد و پایان فریم را مشخص می‌کند. بیت توازن از \لر{XOR} کاهشی هفت بیت داده
به‌دست می‌آید:

\begin{equation}
P = D_0 \oplus D_1 \oplus D_2 \oplus D_3 \oplus D_4 \oplus D_5 \oplus D_6
\end{equation}

در نتیجه، \لر{XOR} همه‌ی بیت‌های داده و بیت توازن صفر و توازن زوج برقرار است. جدول
\رجوع{جدول:قالب‌فریم} اجزای فریم را خلاصه می‌کند.

\شروع{لوح}[ht]
\تنظیم‌ازوسط
\شرح{قالب ده‌بیتی فریم اختصاصی آزمایش}
\شروع{جدول}{|c|c|c|}
\خط‌پر
\سیاه ترتیب & \سیاه نماد & \سیاه مقدار یا نقش \\
\خط‌پر \خط‌پر
۱ & \لر{Start} & صفر \\
۲ & \لر{Parity} & \لر{XOR} هفت بیت داده \\
۳ تا ۹ & $D_0$ تا $D_6$ & داده، از بیت کم‌ارزش \\
۱۰ & \لر{Stop} & یک \\
\خط‌پر
\پایان{جدول}
\برچسب{جدول:قالب‌فریم}
\پایان{لوح}

برای ساعت ۵۰ مگاهرتز و نرخ ۱۱۵۲۰۰ بیت بر ثانیه داریم:

\begin{equation}
\mathrm{BIT\_TICKS} = \frac{50{,}000{,}000}{115{,}200} \simeq 434
\end{equation}

نرخ واقعی حاصل از تقسیم صحیح برابر ۱۱۵۲۰۷٫۳۷۳ بیت بر ثانیه و خطای نسبی آن تقریباً
۰٫۰۰۶۴ درصد است. هر فریم ده‌بیتی ۴۳۴۰ چرخه‌ی ساعت و در فرکانس ۵۰ مگاهرتز تقریباً
۸۶٫۸ میکروثانیه طول می‌کشد. پارامتر \لر{BIT\_TICKS} در هر دو پیمانه مقدار پیش‌فرض ۴۳۴
دارد و در شبیه‌سازی می‌تواند با مقدار کوچک‌تری مانند ۲ جایگزین شود.

\section{معماری سامانه و فایل‌های طراحی}

جدول~\رجوع{جدول:پرونده‌ها} مسئولیت پرونده‌های اصلی را نشان می‌دهد.

\شروع{لوح}[ht]
\تنظیم‌ازوسط
\شرح{پیمانه‌ها و مسئولیت آن‌ها در طرح \لر{UART}}
\شروع{جدول}{|c|p{9.5cm}|}
\خط‌پر
\سیاه پرونده & \سیاه مسئولیت \\
\خط‌پر \خط‌پر
\کد{UARTSender.v} & ذخیره‌ی داده، محاسبه‌ی توازن، تولید فریم سریال و اعلام \لر{busy} \\
\کد{UARTReceiver.v} & هم‌زمان‌سازی \لر{rx}، نمونه‌برداری، بازسازی داده و اعتبارسنجی فریم \\
\کد{UARTTop.v} & اتصال مستقیم \لر{tx} فرستنده به \لر{rx} گیرنده \\
\کد{Tester.v} & اجرای رگرسیون جامع و خودکنترل \\
پرونده‌های \لر{tb\_uart\_*} & بررسی مستقل هر الزام عملکردی و حالت خطا \\
\خط‌پر
\پایان{جدول}
\برچسب{جدول:پرونده‌ها}
\پایان{لوح}

خطوط ۱۸ تا ۳۴ پرونده‌ی \کد{UARTTop.v} اتصال مستقیم حلقه‌ی بازگشت و انتقال یکسان پارامتر
زمانی به دو پیمانه را نشان می‌دهند.

\begin{dsdlcodeblock}
\begin{verbatim}
18      UARTSender #(.BIT_TICKS(BIT_TICKS)) sender (
19          .clk(clk),
20          .rstN(rstN),
21          .new_data(new_data),
22          .send_data(send_data),
23          .tx(tx),
24          .busy(busy)
25      );
26
27      UARTReceiver #(.BIT_TICKS(BIT_TICKS)) receiver (
28          .clk(clk),
29          .rstN(rstN),
30          .rx(tx),
31          .rec_data(rec_data),
32          .rec_new_data(rec_new_data),
33          .correct_data(correct_data)
34      );
\end{verbatim}
\end{dsdlcodeblock}

\section{طراحی فرستنده \لر{UART}}

\subsection{رابط و سیاست پذیرش درخواست}

ورودی‌های فرستنده شامل \لر{clk}، بازنشانی فعال‌پایین \لر{rstN}، پالس یک‌چرخه‌ای
\لر{new\_data} و داده‌ی هفت‌بیتی \لر{send\_data} هستند. خروجی \لر{tx} خط سریال و خروجی
\لر{busy} نشان‌دهنده‌ی فعال‌بودن ارسال است. فقط اگر \لر{new\_data} در لبه‌ی بالارونده‌ی
ساعت و در حالت \لر{IDLE} فعال باشد، درخواست پذیرفته می‌شود. داده و توازن همان لحظه ذخیره
می‌شوند؛ ازاین‌رو تغییر ورودی در میانه‌ی ارسال بر فریم جاری اثر ندارد.

\subsection{ماشین حالت و تولید خروجی}

ماشین حالت فرستنده پنج مرحله‌ی \لر{IDLE}، \لر{START}، \لر{PARITY}، \لر{DATA} و
\لر{STOP} دارد. خطوط ۴۵ تا ۴۸ پرونده‌ی \کد{UARTSender.v} تولید مستقیم خروجی‌ها از حالت
جاری را نشان می‌دهند.

\begin{dsdlcodeblock}
\begin{verbatim}
45      assign busy = state != IDLE;
46      assign tx = state == START  ? 1'b0 :
47                  state == PARITY ? parity_latch :
48                  state == DATA   ? data_latch[bit_index] : 1'b1;
\end{verbatim}
\end{dsdlcodeblock}

خطوط ۵۹ تا ۶۶ همان پرونده نشان می‌دهند که درخواست فقط در حالت بیکار پذیرفته و داده و
توازن در ثبات‌های داخلی ذخیره می‌شوند.

\begin{dsdlcodeblock}
\begin{verbatim}
59                  IDLE: begin
60                      tick_count <= {TICK_WIDTH{1'b0}};
61                      bit_index  <= 3'd0;
62                      if (new_data) begin
63                          data_latch   <= send_data;
64                          parity_latch <= ^send_data;
65                          state        <= START;
66                      end
\end{verbatim}
\end{dsdlcodeblock}

شمارنده‌ی \لر{tick\_count} هر نماد را دقیقاً \لر{BIT\_TICKS} چرخه پایدار نگه می‌دارد.
درخواست جدید هنگام \لر{busy} نادیده گرفته می‌شود و نه صف می‌شود و نه فریم را از ابتدا
آغاز می‌کند.

\section{طراحی گیرنده \لر{UART}}

\subsection{هم‌زمان‌سازی و تشخیص شروع}

ورودی \لر{rx} نسبت به \لر{clk} ناهمگام است. خطوط ۴۹ تا ۵۸ پرونده‌ی
\کد{UARTReceiver.v} هم‌زمان‌ساز دو ثباتی را نشان می‌دهند که ماشین حالت را از ورودی خام
جدا می‌کند.

\begin{dsdlcodeblock}
\begin{verbatim}
49      // Two flip-flops isolate the FSM from an asynchronous serial input.
50      always @(posedge clk or negedge rstN) begin
51          if (!rstN) begin
52              rx_meta <= 1'b1;
53              rx_sync <= 1'b1;
54          end else begin
55              rx_meta <= rx;
56              rx_sync <= rx_meta;
57          end
58      end
\end{verbatim}
\end{dsdlcodeblock}

گیرنده در \لر{IDLE} تا مشاهده‌ی صفر هم‌زمان‌شده منتظر می‌ماند و سپس در میانه‌ی بازه‌ی
بیت شروع، صفرماندن خط را بررسی می‌کند. بازگشت خط به یک پیش از نقطه‌ی میانی باعث رد
شروع کاذب می‌شود.

\subsection{نمونه‌برداری و اعتبارسنجی}

پس از تأیید شروع، گیرنده به‌ترتیب توازن، بیت‌های $D_0$ تا $D_6$ و بیت توقف را با فاصله‌ی
یک \لر{BIT\_TICKS} نمونه‌برداری می‌کند. خطوط ۱۱۷ تا ۱۲۵ پرونده‌ی گیرنده شرط نهایی اعتبار
فریم را نشان می‌دهند.

\begin{dsdlcodeblock}
\begin{verbatim}
117                 STOP: begin
118                     if (tick_count == LAST_TICK) begin
119                         tick_count <= {TICK_WIDTH{1'b0}};
120                         rec_new_data <= 1'b1;
121                         if (rx_sync && ((^data_latch) == parity_latch)) begin
122                             rec_data <= data_latch;
123                             correct_data <= 1'b1;
124                         end
125                         state <= IDLE;
\end{verbatim}
\end{dsdlcodeblock}

برای هر فریم کامل، \لر{rec\_new\_data} دقیقاً یک چرخه فعال می‌شود. \لر{correct\_data}
فقط در همان چرخه و برای فریم معتبر یک است. \لر{rec\_data} در فریم نامعتبر تغییر نمی‌کند؛
پس خطای توازن یا توقف داده‌ی معتبر قبلی را از بین نمی‌برد.

\section{روش آزمون و پوشش الزامات}

همه‌ی آزمون‌گرها ساعت با دوره‌ی ۱۰ نانوثانیه تولید می‌کنند، ابتدا بازنشانی را اعمال می‌کنند،
ورودی‌ها را دور از لبه‌ی نمونه‌برداری تغییر می‌دهند و برای انتظارها کران زمانی محدود دارند.
مقایسه‌ها با \لر{case equality} انجام می‌شوند تا مقادیر ناشناخته معتبر تلقی نشوند.
جدول~\رجوع{جدول:آزمون‌ها} نتیجه‌ی اجرای واقعی هر آزمون را خلاصه می‌کند.

\شروع{لوح}[ht]
\تنظیم‌ازوسط
\شرح{پوشش و نتیجه‌ی آزمون‌های خودکنترل}
\شروع{جدول}{|p{4cm}|p{7.5cm}|c|}
\خط‌پر
\سیاه آزمون & \سیاه هدف & \سیاه نتیجه \\
\خط‌پر \خط‌پر
\لر{tb\_uart\_frame\_format} & ترتیب ده نماد و دوام ۴۳۴ چرخه‌ای هر نماد & \لر{PASS} \\
\لر{tb\_uart\_valid} & دریافت معتبر مقدار \لر{7'b1011001} & \لر{PASS} \\
\لر{tb\_uart\_parity\_error} & رد توازن معکوس‌شده & \لر{PASS} \\
\لر{tb\_uart\_stop\_error} & رد بیت توقف صفر & \لر{PASS} \\
\لر{tb\_uart\_back\_to\_back} & دریافت مرتب دو فریم متوالی & \لر{PASS} \\
\لر{tb\_uart\_busy\_ignore} & نادیده‌گرفتن درخواست هنگام \لر{busy} & \لر{PASS} \\
\لر{tb\_uart\_reset\_idle} & مقادیر خروجی در بازنشانی و بیکاری & \لر{PASS} \\
\لر{tb\_uart\_output\_pulses} & پهنای یک‌چرخه‌ای پالس‌ها & \لر{PASS} \\
\لر{tb\_uart\_all\_values} & همه‌ی ۱۲۸ مقدار هفت‌بیتی & \لر{PASS} \\
\لر{tb\_uart\_false\_start} & رد شروع کاذب و بازیابی & \لر{PASS} \\
\لر{Tester} & رگرسیون جامع همه‌ی سناریوها & \لر{PASS} \\
\خط‌پر
\پایان{جدول}
\برچسب{جدول:آزمون‌ها}
\پایان{لوح}

\section{تحلیل شبیه‌سازی‌ها}

\subsection{دریافت معتبر}

برای داده‌ی \لر{7'b1011001} معادل \لر{0x59}، توازن صفر است. شکل~\رجوع{شکل:فرستنده‌معتبر}
باید خط \لر{tx} را با ترتیب بیت شروع، توازن، هفت بیت داده و بیت توقف و همچنین فعال‌بودن
\لر{busy} در سراسر فریم نشان دهد.

\شروع{شکل}[H]
\centerimg{uart_valid_sender.png}{14cm}
\شرح{شکل موج فرستنده برای فریم معتبر \لر{0x59} و ترتیب دقیق ده نماد سریال}
\برچسب{شکل:فرستنده‌معتبر}
\پایان{شکل}

شکل~\رجوع{شکل:گیرنده‌معتبر} باید بازسازی \لر{rec\_data=7'b1011001} و هم‌زمانی پالس‌های
\لر{rec\_new\_data} و \لر{correct\_data} را نشان دهد.

\شروع{شکل}[H]
\centerimg{uart_valid_receiver.png}{14cm}
\شرح{شکل موج گیرنده برای دریافت معتبر \لر{0x59} و پالس‌های یک‌چرخه‌ای پایان و اعتبار}
\برچسب{شکل:گیرنده‌معتبر}
\پایان{شکل}

\subsection{خطای توازن}

پس از دریافت یک مقدار معتبر، مقدار \لر{0x2D} با توازن معکوس‌شده اعمال می‌شود. شکل
\رجوع{شکل:خطای‌توازن} باید فعال‌شدن پالس پایان، صفرماندن اعتبار و حفظ آخرین داده‌ی معتبر را
نشان دهد.

\شروع{شکل}[H]
\centerimg{uart_parity_error.png}{14cm}
\شرح{رد فریم \لر{0x2D} با بیت توازن نادرست و حفظ آخرین \لر{rec\_data} معتبر}
\برچسب{شکل:خطای‌توازن}
\پایان{شکل}

\subsection{خطای بیت توقف}

در شکل~\رجوع{شکل:خطای‌توقف} داده و توازن صحیح‌اند، اما بیت توقف صفر است؛ بنابراین
\لر{rec\_new\_data} پایان فریم را اعلام می‌کند ولی \لر{correct\_data} صفر می‌ماند.

\شروع{شکل}[H]
\centerimg{uart_stop_error.png}{14cm}
\شرح{رد فریم دارای بیت توقف صفر با وجود داده و توازن صحیح}
\برچسب{شکل:خطای‌توقف}
\پایان{شکل}

\subsection{فریم‌های متوالی}

شکل~\رجوع{شکل:فریم‌های‌متوالی} باید دو مقدار \لر{0x59} و \لر{0x2D} را در نخستین فرصت
قانونی متوالی و دو پالس پایان مرتب نشان دهد.

\شروع{شکل}[H]
\centerimg{uart_back_to_back.png}{14cm}
\شرح{ارسال متوالی \لر{0x59} و \لر{0x2D} و دریافت دو نتیجه‌ی معتبر به همان ترتیب}
\برچسب{شکل:فریم‌های‌متوالی}
\پایان{شکل}

\subsection{درخواست هنگام \لر{busy}}

در شکل~\رجوع{شکل:درخواست‌هنگام‌بیزی} هنگام ارسال \لر{0x59}، درخواست \لر{0x2D} در بازه‌ی
\لر{busy} اعمال می‌شود. فریم نخست نباید تغییر کند و پالس پایان دومی نباید تولید شود.

\شروع{شکل}[H]
\centerimg{uart_busy_ignore.png}{14cm}
\شرح{نادیده‌گرفتن درخواست دوم در زمان \لر{busy} و نبود فریم یا \لر{completion} دوم}
\برچسب{شکل:درخواست‌هنگام‌بیزی}
\پایان{شکل}

\subsection{بازنشانی و حالت بیکار}

شکل~\رجوع{شکل:بازنشانی} باید نشان دهد که با فعال‌شدن \لر{rstN=0}، خروجی \لر{tx} یک،
\لر{busy} صفر، داده‌ی گیرنده صفر و پرچم‌های دریافت غیرفعال هستند.

\شروع{شکل}[H]
\centerimg{uart_reset_idle.png}{14cm}
\شرح{مقادیر تعریف‌شده‌ی خروجی‌ها هنگام بازنشانی فعال‌پایین و حالت بیکار}
\برچسب{شکل:بازنشانی}
\پایان{شکل}

\subsection{پهنای پالس‌های خروجی}

شکل~\رجوع{شکل:پالس‌ها} باید یک‌چرخه‌ای بودن \لر{rec\_new\_data} و هم‌ترازی
\لر{correct\_data} را در فریم معتبر و صفرماندن آن را در فریم نامعتبر نشان دهد.

\شروع{شکل}[H]
\centerimg{uart_output_pulses.png}{14cm}
\شرح{مقایسه‌ی پالس‌های \لر{completion} و \لر{validity} در فریم معتبر و نامعتبر}
\برچسب{شکل:پالس‌ها}
\پایان{شکل}

\subsection{شروع کاذب}

در شکل~\رجوع{شکل:شروع‌کاذب} یک پالس صفر کوتاه‌تر از نیم بیت نباید \لر{completion} ایجاد کند.
دریافت صحیح فریم بعدی نیز بازیابی گیرنده را اثبات می‌کند.

\شروع{شکل}[H]
\centerimg{uart_false_start.png}{14cm}
\شرح{رد پالس شروع کوتاه و دریافت صحیح فریم بعدی پس از بازیابی}
\برچسب{شکل:شروع‌کاذب}
\پایان{شکل}

\subsection{بررسی همه‌ی مقادیر هفت‌بیتی}

آزمون جامع، مقادیر \لر{0x00} تا \لر{0x7F} را به‌ترتیب ارسال کرد و برای هر مقدار دقیقاً یک
پایان معتبر با داده‌ی برابر ورودی به‌دست آمد. شکل‌های \رجوع{شکل:همه‌مقادیر۱} تا
\رجوع{شکل:همه‌مقادیر۴} چهار بخش این پیمایش را نمایش می‌دهند.

\شروع{شکل}[H]
\centerimg{uart_all_values_1.png}{14cm}
\شرح{بخش نخست پیمایش جامع مقادیر هفت‌بیتی از ابتدای دامنه}
\برچسب{شکل:همه‌مقادیر۱}
\پایان{شکل}

شکل~\رجوع{شکل:همه‌مقادیر۲} ادامه‌ی پیمایش مقادیر هفت‌بیتی را نشان می‌دهد.

\شروع{شکل}[H]
\centerimg{uart_all_values_2.png}{14cm}
\شرح{بخش دوم پیمایش جامع مقادیر هفت‌بیتی}
\برچسب{شکل:همه‌مقادیر۲}
\پایان{شکل}

شکل~\رجوع{شکل:همه‌مقادیر۳} بخش سوم پیمایش را نمایش می‌دهد.

\شروع{شکل}[H]
\centerimg{uart_all_values_3.png}{14cm}
\شرح{بخش سوم پیمایش جامع مقادیر هفت‌بیتی}
\برچسب{شکل:همه‌مقادیر۳}
\پایان{شکل}

شکل~\رجوع{شکل:همه‌مقادیر۴} پایان پیمایش تا \لر{0x7F} را نشان می‌دهد.

\شروع{شکل}[H]
\centerimg{uart_all_values_4.png}{14cm}
\شرح{بخش پایانی پیمایش تا \لر{0x7F} و دریافت معتبر آخرین مقدار}
\برچسب{شکل:همه‌مقادیر۴}
\پایان{شکل}

\section{خروجی واقعی رگرسیون جامع}

کد موجود با \لر{Icarus Verilog 12.0} کامپایل و اجرا شد. خروجی واقعی پیمانه‌ی
\لر{Tester} به‌صورت زیر بود:

\begin{dsdlcodeblock}
\begin{verbatim}
PASS: reset and idle outputs
TEST: exact valid frame and symbol timing
TEST: request asserted while busy
TEST: required consecutive frames
TEST: exhaustive seven-bit loopback
TEST: valid raw receiver frame
TEST: required parity error
TEST: required stop-bit error
TEST: false-start rejection and recovery
PASS: complete DSDL7 UART regression
\end{verbatim}
\end{dsdlcodeblock}

هر ده آزمون مستقل نیز با کد خروج صفر و پیام \لر{PASS} پایان یافت. آزمون قالب فریم دوام
۴۳۴ چرخه‌ای هر نماد، آزمون جامع تمام ۱۲۸ داده و آزمون‌های خطا رفتار مورد انتظار را تأیید
کردند.

\section{بررسی سنتز}

طرح با \لر{Yosys 0.33} خوانده و به سلسله‌مراتب \لر{UARTTop} تبدیل شد. گذر
\لر{check} صفر مشکل گزارش کرد. پس از تبدیل فرایندها و بهینه‌سازی، گیرنده ۶۳ سلول، فرستنده
۴۵ سلول و کل سلسله‌مراتب ۱۰۸ سلول داشت. این اعداد نتیجه‌ی بررسی عمومی \لر{RTL} هستند و
جایگزین گزارش وابسته به قطعه‌ی هدف در \لر{Quartus} نیستند.

شکل~\رجوع{شکل:سنتزکواتروس} باید پایان موفق مرحله‌ی \لر{Analysis \& Synthesis} را نشان دهد.

\شروع{شکل}[H]
\centerimg{quartus_analysis_synthesis.png}{14cm}
\شرح{نتیجه‌ی مرحله‌ی \لر{Analysis \& Synthesis} در \لر{Quartus} و نبود خطای طراحی}
\برچسب{شکل:سنتزکواتروس}
\پایان{شکل}

شکل~\رجوع{شکل:نمای‌آرتی‌ال} باید سلسله‌مراتب \لر{UARTTop} و اتصال مستقیم فرستنده و گیرنده
را نمایش دهد.

\شروع{شکل}[H]
\centerimg{quartus_rtl_viewer.png}{14cm}
\شرح{نمای \لر{RTL} شامل \لر{UARTTop}، \لر{UARTSender}، \لر{UARTReceiver} و اتصال مستقیم \لر{tx} به \لر{rx}}
\برچسب{شکل:نمای‌آرتی‌ال}
\پایان{شکل}

شکل~\رجوع{شکل:منابع‌کواتروس} باید خلاصه‌ی منابع مصرفی قطعه‌ی انتخاب‌شده را نشان دهد.

\شروع{شکل}[H]
\centerimg{quartus_resource_utilization.png}{14cm}
\شرح{خلاصه‌ی \لر{Resource Utilization} پروژه‌ی \لر{Quartus} برای قطعه‌ی هدف}
\برچسب{شکل:منابع‌کواتروس}
\پایان{شکل}

\section{نتیجه‌گیری}

فرستنده و گیرنده‌ی \لر{UART} هفت‌بیتی با قالب اختصاصی \لر{parity-before-data} پیاده‌سازی
شدند. فرستنده با ذخیره‌ی داده در لحظه‌ی پذیرش درخواست و نادیده‌گرفتن درخواست‌های زمان
\لر{busy} از تغییر فریم فعال جلوگیری می‌کند. گیرنده با هم‌زمان‌ساز دو مرحله‌ای، تأیید شروع
در میانه‌ی بیت و نمونه‌برداری هم‌فاز، داده را بازسازی و خطاهای توازن و توقف را تشخیص می‌دهد.
یک‌چرخه‌ای بودن \لر{rec\_new\_data} و \لر{correct\_data} امکان تشخیص فریم‌های متوالی را
فراهم می‌سازد. موفقیت یازده شبیه‌سازی و بررسی صفرخطای \لر{Yosys} نشان می‌دهد \لر{RTL}
موجود الزامات عملکردی و ساختاری تعریف‌شده را برآورده می‌کند. تصاویر شکل موج و سه خروجی
\لر{Quartus} باید پیش از نسخه‌ی نهایی قابل ارائه در محل‌های تعیین‌شده قرار گیرند.
